\section{The unit ball of theta expansions and its primitive frame}
\label{sec:theta-unit-basis}

Keep the valued algebraic closure \(F\), its valuation ring
\(\mathcal V\), and the PEL degeneration over the complete discrete
valuation ring \(T\). Let \(B/T\) be the abelian part of dimension
\(d\), with its principal theta bundle \(L\), and let \(X\) be the
period lattice of rank \(r\). At level \(q\ge1\) the coefficient for
\(\chi\in X\) is a section of a line bundle \(L_{q,\chi}\) on \(B\).
In a rigidification its class is \(L^{\otimes q}\) times the
degree-zero twist attached to \(\chi\). We compute the unit ball of
the expansion space in the coefficient valuation of
Section~\ref{sec:theta-lift-valuation}.

\paragraph{Integral coordinates at one character.}
The module \(V_{q,\chi}=H^0(B,L_{q,\chi})\) is finite over \(T\):
the base is noetherian, \(B\) proper, and the line bundle coherent.
It is torsion free, because \(B\) is flat over \(T\) and a line
bundle is locally free. Thus \(V_{q,\chi}\) is free over \(T\).
Flat base change identifies \(V_{q,\chi}\otimes_T F\) with the
sections on \(B_F\). On that fiber the bundle is ample, since its
degree-zero twist preserves numerical equivalence. Its polarization
homomorphism is \(q\lambda_L\), of degree \(q^{2d}\) by the principal
polarization. Riemann--Roch for abelian varieties gives
\(\chi(L_{q,\chi})^2=q^{2d}\); ampleness gives positive Euler
characteristic and vanishing of higher cohomology. Hence
\(\operatorname{rank}_T V_{q,\chi}=q^d\).
\footnote{\href{https://stacks.math.columbia.edu/tag/02O6}{Stacks Project,
Lemma~30.19.2, proper cohomology over an affine noetherian base};
\href{https://stacks.math.columbia.edu/tag/02KH}{Lemma~30.5.2, flat base change};
\href{https://jmilne.org/math/CourseNotes/AV.pdf}{J.~S.~Milne,
\emph{Abelian Varieties}, Chapter~I, Theorems~7.2 and 11.1};
\href{https://wiki.epfl.ch/shimvar/documents/IHESlectures-Genestier-Ngo.pdf}{A.~Genestier
and B.~C.~Ng\^o, \emph{Lectures on Shimura varieties},
Theorems~2.2.1--2.2.2, Riemann--Roch and the index of an ample class}.}

Choose a \(T\)-basis \(b_{\chi,1},\ldots,b_{\chi,q^d}\).
Writing a coefficient as \(\sum_i c_i b_{\chi,i}\) over \(F\),
its section valuation is at least \(h\in\mathbb Q\) if and only
if \(v(c_i)\ge h\) for each \(i\). The reverse direction follows
by taking this basis as the integral sections in the definition.
For the forward direction, expand the integral sections in any
given representation in this basis. Their coordinates lie in \(T\),
so each resulting coefficient is a sum of terms of valuation at
least \(h\). This interpretation gives all rational thresholds,
including when \(F\) is nondiscretely valued.

\paragraph{Coset coordinates for the whole family.}
Choose representatives \(e_j\) for \(X/qX\), of cardinality \(q^r\).
A theta expansion is determined by its coefficients at the \(e_j\).
Indeed, its functional equation determines the coefficient at
\(e_j+qz\) by translating that at \(e_j\) by \(c(z)\) on the
abelian part and tensoring with the period multiplier. Each
\(\chi\) has one such expression after fixing \(e_j\), because \(X\)
has no torsion. Conversely assign an arbitrary coefficient at each
representative and use this rule for all \(z\). Applying the rule
twice agrees with its application to the sum of the periods by
the period cocycle. This constructs a family and proves that
the joint coordinate map is a linear isomorphism
\[
     s\longmapsto(c_{j,i}(s))_{j,i},\qquad
     \mathcal F_q \ \simeq\
               \bigoplus_{j=1}^{q^r} F^{q^d}.
\]
Here \(\mathcal F_q\) denotes the space of expansion families
with the given cubical descent law.

In \(\chi=e_j+qz\) the shift in the lower bound of the coefficient
is \(qQ(z)+B(z,e_j)\). The identity
\[
 qQ(z)+B(z,e_j)
       =\frac{Q(e_j+qz)-Q(e_j)}{q}
\]
shows that the normalized bound at \(e_j\) propagates to its
whole coset. Thus the full filtration is given in these coordinates by
\[
 \mathcal U_h(s)\quad\Longleftrightarrow\quad
          v(c_{j,i}(s))\ge h+Q(e_j)/q\quad\hbox{for all }j,i .
\]
For each \(j\), choose \(a_j\in F^\times\) with
\(a_j^q=\pi^{Q(e_j)}\), possible as \(F\) is algebraically closed.
Lift the coordinate vector whose only nonzero entry is \(a_j\) in
position \((j,i)\) to its family. The resulting \(q^{r+d}\) families
are independent over \(F\) and freely generate \(\mathcal U_0\)
over \(\mathcal V\). Division by \(a_j\) gives the integral
coordinates of any element of the unit ball.
\footnote{\href{https://github.com/ChristianKudeba/alpoge-rh-verification/blob/alpoge-source-7cbf7a0c/src/Development/AbelianVarietyTheory/Proofs/RingsOfIntegers/ExistsBasisOfEq.lean}{Valuation in integral coefficient bases};
\href{https://github.com/ChristianKudeba/alpoge-rh-verification/blob/alpoge-source-7cbf7a0c/src/ReciprocityTheory/ReductiveGroups/ArithmeticGeometry/DiscreteValuationRings/ExistsFamRep.lean}{extension from one coset};
\href{https://github.com/ChristianKudeba/alpoge-rh-verification/blob/alpoge-source-7cbf7a0c/src/Development/AbelianVarietyTheory/Proofs/RingsOfIntegers/ExistsBasisOfCoord.lean}{scaled coordinate basis of the unit ball}.}

\paragraph{The level-one frame and its pullback.}
At \(q=1\) there is one coset with representative zero, and its
coefficient module is free of rank one. Choose a \(T\)-generator
\(\sigma\); its reduction is nonzero. Indeed the exact sequence
\(
 0\longrightarrow L_{1,0}\xrightarrow{\pi}L_{1,0}
 \longrightarrow L_{1,0}|_{B_{\kappa(T)}}\longrightarrow0
\)
identifies the kernel of reduction on sections with
\(\pi V_{1,0}\). Extend \(\sigma\) from zero by the period
functional equation to \(\theta\). As \(Q(0)=0\),
\(\theta\in\mathcal U_0\) at level one, and its zero coefficient
is the base change of a primitive integral section. Positivity of
\(Q\) also implies that in the base-field integral expansion all
coefficients outside zero are divisible by \(\pi\).

Let \(\alpha\overline\alpha=q\). In the level-\(q\) pullback the
coefficient at \(\overline\alpha\chi\) is the pullback of the
\(\chi\)-coefficient by \([\alpha]\), using the normalized
identification of its line; all other coefficients vanish.
Here the choice of theta bundle in the PEL uniformization matters.
Write \(L_x\) for its pullback to a point \(x\),
\(\ell(x)=L_x\otimes L_0^{-1}\), and
\[
 \mathcal P(x,y)=L_{x+y}\otimes L_0\otimes L_x^{-1}\otimes L_y^{-1}.
\]
The supplied integral isomorphism for the distinguished element \(w_0\)
identifies \(\ell(x)\) with \(\mathcal P(x,w_0x)\), respecting the
rigidification at zero. This is the extra line-bundle datum behind the
pullback formula. By the Rosati adjoint condition on the biextension,
bilinearity, and commutation of the CM endomorphisms, there are
rigidified isomorphisms over \(T\)
\[
 \ell(\alpha x)\simeq\mathcal P(\alpha x,w_0\alpha x)
  \simeq\mathcal P(x,\overline\alpha w_0\alpha x)
  =\mathcal P(x,qw_0x)\simeq\ell(x)^{\otimes q}.
\]
The coefficient line at \(x\) for \((k,\chi)\) is
\(L_x^{\otimes k}\otimes\mathcal P(x,c(\chi))\).
The same adjoint identity sends its biextension factor under
\(\alpha\) to \(\mathcal P(x,\overline\alpha c(\chi))
=\mathcal P(x,c(\overline\alpha\chi))\). Taking the preceding
isomorphism to the \(k\)-th tensor power and inserting the
rigidification of \(L_0^{\otimes k(q-1)}\) therefore identifies
the pulled-back coefficient line with \(L_{kq,\overline\alpha\chi}\).
The generic generator in the normalized pullback agrees with this
integral identification: their quotient is a global unit on the
proper geometrically connected abelian part over the generic field,
hence a scalar, and their common normalization at the identity sets
it equal to one. This proves the asserted absence of a scalar
valuation loss. The step uses the integral \(w_0\)-isomorphism
as well as Rosati compatibility.
\footnote{\href{https://github.com/ChristianKudeba/alpoge-rh-verification/blob/alpoge-source-7cbf7a0c/src/ReciprocityTheory/ReductiveGroups/ArithmeticGeometry/DiscreteValuationRings/PullGenPf.lean}{Integral square identity, character twist, and comparison of normalized generators}.}
The balanced CM action gives
\[
    2Q(\overline\alpha\chi)
       =B(\overline\alpha\chi,\overline\alpha\chi)
       =B(\chi,\alpha\overline\alpha\chi)
       =2qQ(\chi).
\]
Thus the required normalized bounds agree. Pullback of integral
sections preserves their integrality and gives
\(\theta_\alpha\in\mathcal U_0\) at level \(q\).
\footnote{\href{https://github.com/ChristianKudeba/alpoge-rh-verification/blob/alpoge-source-7cbf7a0c/src/Development/AbelianVarietyTheory/Proofs/JetHarmonics/ExistsThetaFrame.lean}{Primitive level-one frame};
\href{https://github.com/ChristianKudeba/alpoge-rh-verification/blob/alpoge-source-7cbf7a0c/src/ReciprocityTheory/ReductiveGroups/ArithmeticGeometry/DiscreteValuationRings/QuadraticAlgebraValuation.lean}{CM scaling of \(Q\)};
\href{https://github.com/ChristianKudeba/alpoge-rh-verification/blob/alpoge-source-7cbf7a0c/src/ReciprocityTheory/ReductiveGroups/ArithmeticGeometry/DiscreteValuationRings/PullbackValuationInvarianceOf.lean}{forward pullback on the coefficient bounds}.}

To exclude positive depth, suppose \(\mathcal U_h(\theta_\alpha)\)
holds for some \(h>0\). At the zero index, \(Q(0)=0\). Therefore
this coefficient is an integral combination with scalars in the
maximal ideal of \(\mathcal V\), and it vanishes upon reduction
to the residue field of \(\mathcal V\). The normalized pullback
identification at zero is an isomorphism of the integral lines,
so \([\alpha]^*\overline\sigma=0\) on this abelian special fiber.
Further pullback by \([\overline\alpha]\) gives
\([q]^*\overline\sigma=0\). Multiplication by \(q\) on an abelian
variety is finite faithfully flat, so pullback detects a zero
section. But \(\overline\sigma\ne0\): the original reduction was
nonzero and the extension of residue fields is faithfully flat.
This contradiction proves primitivity of the pullback for the
coefficient valuation. Together with
Section~\ref{sec:theta-lift-valuation}, it gives primitivity of
every normalized finite-order translate used in the span argument.
\footnote{\href{https://github.com/ChristianKudeba/alpoge-rh-verification/blob/alpoge-source-7cbf7a0c/src/ReciprocityTheory/ReductiveGroups/ArithmeticGeometry/FirstCohomology/LevelFrameOfLaws.lean}{Pullback of the zero coefficient and descent of divisibility through multiplication by \(q\)}.}

This calculation assumes the PEL degeneration on the given chart,
its rigidified integral \(w_0\) and Rosati comparisons, and the
cocycle of its period multipliers. The scalar and rank calculations
are valid in either residue characteristic once these data are
given. The use of an integral eigenbasis of torsion translates in
Section~\ref{sec:theta-translate-span} additionally requires
\(q\) to be a unit.
